\section{Operational Workflow}
\label{sec:workflow}
This section describes how support teams deploy Afar-prgenv and how users
interact with it during early drop evaluation.

\begin{figure*}[t]
  \centering
  \begin{tikzpicture}[
  wfstep/.style={draw, rounded corners, align=center, minimum width=2.4cm, minimum height=0.9cm},
  arrow/.style={-Latex, thick}
]
  \node[wfstep] (drop) {AFAR\\drop};
  \node[wfstep, right=1.2cm of drop] (gen) {Generate\\modules};
  \node[wfstep, right=1.2cm of gen] (repro) {Repro\\tests};
  \node[wfstep, right=1.2cm of repro] (validate) {Validation\\suites};
  \node[wfstep, right=1.2cm of validate] (ready) {User-ready\\release};

  \draw[arrow] (drop) -- (gen);
  \draw[arrow] (gen) -- (repro);
  \draw[arrow] (repro) -- (validate);
  \draw[arrow] (validate) -- (ready);
\end{tikzpicture}

  \caption{Monthly AFAR drop workflow. Afar-prgenv provides a repeatable path
  from drop ingestion to user readiness.}
  \Description{Vertical workflow from AFAR drop ingestion through module
  generation, repro tests, validation suites, and user-ready release.}
  \label{fig:workflow}
\end{figure*}

\subsection{Support Team Workflow}
\label{sec:support-workflow}
Support staff add a new AFAR drop under the configured root, update the version
mapping, and regenerate the module tree using
\texttt{afar\_modules/scripts/generate\_afar\_modules.sh}. They then run the repro test suite
via \texttt{scripts/run\_afar\_tests.sh --keep-going} and capture logs for validation.
A drop is considered ready for user evaluation when the repro matrix passes
and the wrapper diagnostics show stable MPI and library metadata.
The harness writes per-profile logs under \texttt{logs/}, which serve as
support artifacts for regression tracking.

When PE library modules do not provide usable `pkg-config` metadata, the
support workflow also includes generating shim files via
\texttt{generate\_pkgconfig\_shims.sh}. This step can be
triggered explicitly or left to the wrapper auto-refresh logic, but it is
recorded in logs so that metadata changes are traceable during drop triage.

After repro gating, support staff run extended validation suites, including an
OLCF Fortran application set and the SOLLVE validation suite, to assess
production readiness beyond basic compile and link coverage.

\subsection{Admin Deployment}
\label{sec:admin}
For site deployment, the AFAR module tree (\texttt{modulefiles}, \texttt{bin},
\texttt{pkgconfig}, and configuration files) is copied into a site modulepath.
Admins ensure wrapper scripts are executable and add the module path via
`module use`. Updates to AFAR drops or MPICH paths are performed by refreshing
the version mapping and re-running the generator, followed by a smoke test
(`ftn`, `pkg-config --cflags mpichf90`) to confirm wrapper and metadata health.

\subsection{User Workflow}
\label{sec:user-workflow}
Users follow the standard PE module sequence (e.g., `cpe/\peversion`,
`PrgEnv-amd/\prgenvversion`) and then load `afar-prgenv/\afarversion`. The
wrappers preserve the expected compile and link commands while injecting AFAR
paths and offload flags when needed.

\subsection{Failure Modes and Diagnostics}
\label{sec:diagnostics}
Common failure modes include missing or mismatched MPI metadata, incompatible
Fortran module files from vendor libraries, and absent ROCm runtime libraries.
Wrapper diagnostics surface these issues early by reporting filtered flags,
active MPI flavors, and updated \texttt{PKG\_CONFIG\_PATH} settings.

The repro harness also exercises different module load orders to ensure that
the wrappers re-synchronize MPI and library metadata correctly when users load
modules before or after `afar-prgenv`.
